\documentclass[11pt]{article}
\usepackage[T1]{fontenc}
\usepackage[a4paper,margin=1in]{geometry}
\usepackage{amsmath,amssymb,amsthm}
\usepackage[hidelinks]{hyperref}
\usepackage{microtype}

\newtheorem*{question}{Source open problem}
\newtheorem*{answer}{Later theorem (specialized to codimension one)}

\setlength{\parindent}{0pt}
\setlength{\parskip}{0.6em}

\title{Literature Status: Energy-Concentrated Wave-Packet\\
Frames on Higher-Dimensional Euclidean Balls}
\author{Run \texttt{fa\_banach\_001} --- lane 0}
\date{9 August 2026}

\begin{document}
\maketitle

\begin{center}
\fbox{\parbox{0.9\textwidth}{
\textbf{Status: literature already answered.}
The higher-dimensional problem recorded after Theorem~1.2 of
arXiv:2601.21224 is answered by the same authors in the separate later paper
arXiv:2607.23953.  Its Theorem~1.2 constructs the requested frame for every
$d\geq2$ and obtains the desired $R^{d-1}$ bound up to a polylogarithmic
factor.
}}
\end{center}

\section{Original question}

K. Hughes, A. Israel, and A. Mayeli,
\emph{Wave Packets and Eigenvalue Estimates for Limiting Operators on the
Disk}, arXiv:2601.21224~\cite{disk}, proves the relevant energy-concentration
theorem for $d=2$.  Immediately after Theorem~1.2, on arXiv PDF page~3, the
authors identify the remaining case as follows.

\begin{question}
For $d\geq3$, construct a wave-packet frame on the $d$-dimensional Euclidean
ball whose index set decomposes as
$\mathcal I=\mathcal I_1\cup\mathcal I_2\cup\mathcal I_3$, with
\[
 \sum_{\nu\in\mathcal I_1}\|\widehat\psi_\nu\|_{L^2(S)}^2
 +\sum_{\nu\in\mathcal I_2}
   \|\widehat\psi_\nu\|_{L^2(\mathbb R^d\setminus S)}^2
 \leq \varepsilon^2,
\]
and with residual cardinality
\[
 \#\mathcal I_3\leq C R^{d-1}\log(R/\varepsilon)^J
\]
for some $J>0$ and $C>0$.
\end{question}

The source states explicitly that its theorem answers this problem for
$d=2$ and that dimensions $d\geq3$ remain open.

\section{Separate later answer}

The same authors subsequently wrote
\emph{Localized frames on Euclidean balls}, arXiv:2607.23953~\cite{balls}.
Its introduction explicitly describes the paper as extending the disk
framework to Euclidean balls in arbitrary dimensions.  Theorem~1.2, on arXiv
PDF pages~2--3, supplies the exact requested construction.

For the codimension-one case $\eta=1$, specialize the theorem's geometric
error term to obtain the following.

\begin{answer}
Fix $d\geq2$ and $s>1$.  For every dyadic $R\geq2$ there is a unit-norm frame
$\{\psi_\nu\}_{\nu\in\mathcal I}$ for $L^2(B_d(R))$, with frame bounds
depending only on $d$.  For every eligible bounded frequency set $S$ and
$\varepsilon\in(0,1/2]$, its index set has the energy decomposition in the
source question and
\[
 \#\mathcal I_3
 \lesssim_{d,s}
 \mathcal M^{d-1}(\partial S)R^{d-1}
 \log\!\left(\frac{R\operatorname{diam}(S)}{\varepsilon}\right)^{sd+1}.
\]
\end{answer}

After the standard normalization of $\operatorname{diam}(S)$, this is the
required estimate with $J=sd+1$.  The dyadic restriction is removed by the
rescaling noted directly after the theorem.  Thus the later paper answers all
dimensions $d\geq3$ from the source problem.  It is stronger in another
direction as well: its full statement treats boundaries of fractional
codimension $\eta\in(0,1]$.

\section{Status and bounded search}

A bounded search through 9 August 2026 checked the run registry and cheap
local indexes for arXiv:2601.21224, followed by arXiv searches using the exact
title, authors, and the phrases ``wave packet frame'', ``Euclidean ball'', and
``higher dimensions''.  The decisive match was arXiv:2607.23953.  Its theorem
and its explicit relation to the disk paper were checked in both the arXiv
source and PDF.  This packet records an already-known later answer and claims
no new mathematical result.

\begin{thebibliography}{9}
\bibitem{disk}
K. Hughes, A. Israel, and A. Mayeli,
\emph{Wave Packets and Eigenvalue Estimates for Limiting Operators on the
Disk}, arXiv:2601.21224 (2026).

\bibitem{balls}
K. Hughes, A. Israel, and A. Mayeli,
\emph{Localized frames on Euclidean balls}, arXiv:2607.23953 (2026).
\end{thebibliography}

\end{document}
